\documentclass[a4paper]{article}
\usepackage[utf8]{inputenc}
\usepackage[T1]{fontenc}
\usepackage{textcomp}
\usepackage[english]{babel}
\usepackage{amsmath, amssymb, amsthm}
\usepackage{hyperref}
\usepackage{cleveref}
\usepackage{geometry}

% figure support
\usepackage{import}
\usepackage{xifthen}
\pdfminorversion=7
\usepackage{pdfpages}
\usepackage{transparent}
\newcommand{\incfig}[1]{%
  \def\svgwidth{\columnwidth}
  \import{./Figures/}{#1.pdf_tex}
}

\usepackage[most]{tcolorbox}

\pdfsuppresswarningpagegroup=1

\newcommand{\R}{\mathbb{R}}
\newcommand{\Q}{\mathbb{Q}}
\newcommand{\Z}{\mathbb{Z}}
\newcommand{\N}{\mathbb{N}}
\renewcommand{\H}{\mathbb{H}}
\newcommand{\F}{\mathbb{F}}
\newcommand{\W}{\mathbb{W}}

\usepackage{tikz}
\usepackage{xparse}

% Define a new environment for polar grid
\NewDocumentEnvironment{polargrid}{O{4}O{4}O{12}O{4.5}}{%
  % #1 = max radius (default: 4)
  % #2 = number of concentric circles (default: 4)
  % #3 = number of rays (default: 12)
  % #4 = axis length beyond max radius (default: 4.5)
  \begin{tikzpicture}[>=latex]
    % Draw the outer circle
    \draw[lightgray, thin] (0,0) circle (#1);
    % Draw concentric circles
    \foreach \r in {1,...,\numexpr#2-1\relax} {
      \draw[lightgray, thin] (0,0) circle (\r*#1/#2);
    }
    % Draw rays
    \foreach \i in {0,...,\numexpr#3-1\relax} {
      \pgfmathsetmacro{\angle}{(\i*360)/#3}
      \draw[lightgray, thin] (\angle:#1) -- (\angle:0);
    }
    % Draw main axes
    \foreach \angle in {0,90,180,270} {
      \draw[black, thick] (\angle:#1) -- (\angle:0);
    }
    % Angle labels
    \foreach \i/\label in {
      30/$\frac{\pi}{6}$,
      60/$\frac{\pi}{3}$,
      120/$\frac{2\pi}{3}$,
      150/$\frac{5\pi}{6}$,
      180/$\pi$,
      210/$\frac{7\pi}{6}$,
      240/$\frac{4\pi}{3}$,
      270/$\frac{3\pi}{2}$,
      300/$\frac{5\pi}{3}$,
      330/$\frac{11\pi}{6}$
    } {
      \draw (\i:#1 + 0.3) node {\label};
    }
    % Draw axis labels
    \draw[->, thick] (0,0) -- (0,#4) node[above] {$\mathbf{y}$};
    \draw[->, thick] (0,0) -- (#4,0) node[right] {$\mathbf{x}$};
  }{%
  \end{tikzpicture}
}

% Command for plotting points on polar grid
% This version properly handles negative radii by adding 180 degrees to the angle
\NewDocumentCommand{\polarpoint}{mmO{red}O{}}{%
  % #1 = radius
  % #2 = angle in radians (will be converted to degrees)
  % #3 = point color (default: red)
  % #4 = label (optional)

  % Handle negative radius by adding 180° to angle and using absolute radius
  \pgfmathsetmacro{\radius}{abs(#1)}
  \pgfmathsetmacro{\angle}{#2*180/pi}

  % If radius is negative, adjust the angle by adding 180°
  \ifdim #1 pt < 0 pt
    \pgfmathsetmacro{\angle}{\angle + 180}
  \fi

  % Normalize angle to be between 0° and 360°
  \pgfmathsetmacro{\angle}{mod(\angle, 360)}

  \draw[fill=#3] (\angle:\radius) circle (2pt) 
  \ifx\relax#4\relax\else
    node[above right, font=\scriptsize] {#4}
  \fi;
}

% For convenience, a command that formats polar coordinates nicely
\NewDocumentCommand{\polarcoord}{mm}{%
  % #1 = radius
  % #2 = angle (in radians)
  $(#1, \formatangle{#2})$%
}

% Helper command to format angles nicely
\NewDocumentCommand{\formatangle}{m}{%
  \ensuremath{#1}%
}

\usepackage[most]{tcolorbox}
\tcbuselibrary{listings, breakable, skins}

% Problem box
\newtcolorbox[auto counter, number within=section]{problem}[1][]{%
  colback=red!5!white,
  colframe=red!100!black,
  fonttitle=\bfseries,
  title=Problem~\thetcbcounter,
  label=#1,
  enhanced,
  breakable,
  sharp corners,
  boxrule=0.8pt,
  coltitle=black,
  before upper={\parindent15pt}, % optional formatting
}

% Solution box
\newtcolorbox[auto counter, number within=section]{solution}[1][]{%
  colback=green!5!white,
  colframe=green!100!black,
  fonttitle=\bfseries,
  title=Solution~\thetcbcounter,
  label=#1,
  enhanced,
  breakable,
  sharp corners,
  boxrule=0.8pt,
  coltitle=black,
  before upper={\parindent15pt}, % optional formatting
}

\author{Bhashkar Paudyal}
\title{Workbook 1}

\begin{document}
  \maketitle
  \newpage

  \tableofcontents{}
  \newpage

  \section{Deceptive Coordinates}
  \label{sec:deceptive-coordinates}

  \begin{problem}{Find 5 different ways to express the following using polar
    coordinates.\\}
    \begin{polargrid}[5][5][24][5.5]
      \polarpoint{3}{7*pi/6}[red][\((3,\frac{7\pi}{6})\)]
    \end{polargrid}
    \begin{solution}
      \begin{enumerate}
        \item \((3,\frac{7\pi}{6})\)
        \item \(-3,\frac{\pi}{6}\) 
        \item \((3,-\frac{5\pi}{6})\) 
        \item \((-3,\frac{13\pi}{6})\) 
        \item \((3, \frac{19\pi}{6})\)
      \end{enumerate}
    \end{solution}
  \end{problem}

  \newpage

  \section{Plotting Polar Coordinates}
  \label{sec:plotting-polar-coordinates}

  \textbf{Plot the following polar coordinates.\\}
    \begin{enumerate}
      \item \((5,\frac{3\pi}{4})\) 
      \item \((-4,\frac{17\pi}{3})\) 
      \item \((-4,\frac{\pi}{6})\) 
      \item \((3,-\frac{2\pi}{3})\) 
      \item \((6,-\frac{37\pi}{3})\)
      \item \((-6,-\frac{7\pi}{6})\)
    \end{enumerate}
    
    \begin{polargrid}[6][6][12][6.5]
      \polarpoint{-4}{pi/6}[red][\((-4,\frac{\pi}{6})\)]
      \polarpoint{5}{3*pi/4}[red][\((5,\frac{3\pi}{4})\)]
      \polarpoint{3}{-2*pi/3}[red][\((3,-\frac{2\pi}{3})\)]
      \polarpoint{-6}{-7*pi/6}[red][\((-6,-\frac{7\pi}{6})\)]
      \polarpoint{-4}{17*pi/3}[red][\((-4,\frac{17\pi}{3})\)]
      \polarpoint{6}{-37*pi/3}[red][\((6,-\frac{37\pi}{3})\)]
    \end{polargrid}

    \newpage
    \section{Transformation of Coordinates}
    \label{sec:transformation-of-coordinates}
    We know:
      \begin{equation}
        \label{eqn:polar-pythagorean}
        \(r = \sqrt{x^2+y^2} \)
      \end{equation}
      \begin{equation}
        \label{eqn:polar-theta}
        \(\theta = \tan^{-1}(\frac{y}{x})\)
      \end{equation}
      \begin{equation}
        \label{eqn:cart-x}
        \(x = r\cos\theta\)
      \end{equation}
      \begin{equation}
        \label{eqn:cart-y}
        \(y = r\sin\theta\)
      \end{equation}

\subsection*{Convert to Polar Form}

\begin{enumerate}
    \item \((-5, -5)\)
    
    \textbf{Solution:}
    
    Take \(x = -5\), \(y = -5\). From \crefrange{eqn:polar-pythagorean}{eqn:polar-theta},
    \[
    r = \sqrt{(-5)^2 + (-5)^2} = \sqrt{50} = 5\sqrt{2}
    \]
    \(\theta\) lies in the \(3^{rd}\) quadrant. So,
    \[
    \theta = \tan^{-1}\left(\frac{-5}{-5}\right) + \pi = \pi + \frac{\pi}{4} = \frac{5\pi}{4}
    \]
    
    \(\therefore\) The required point is \(_p(5\sqrt{2}, \frac{5\pi}{4})\).

    \item \(\left(4\sqrt{3}, 4\right)\)
    
    \textbf{Solution:}
    
    Take \(x = 4\sqrt{3}\), \(y = 4\). From \crefrange{eqn:polar-pythagorean}{eqn:polar-theta},
    \[
    r = \sqrt{(4\sqrt{3})^2 + 4^2} = \sqrt{48 + 16} = 8
    \]
    \(\theta\) lies in the \(1^{st}\) quadrant. So,
    \[
    \theta = \tan^{-1}\left(\frac{4}{4\sqrt{3}}\right) = \frac{\pi}{6}
    \]
    
    \(\therefore\) The required point is \(_p(8, \frac{\pi}{6})\).

    \item \((0, -3)\)
    
    \textbf{Solution:}
    
    Take \(x = 0\), \(y = -3\). From \crefrange{eqn:polar-pythagorean}{eqn:polar-theta},
    \[
    r = \sqrt{0^2 + (-3)^2} = 3
    \]
    \(\theta\) lies on the negative \(y\)-axis. So,
    \[
    \theta = \frac{3\pi}{2}
    \]
    
    \(\therefore\) The required point is \(_p(3, \frac{3\pi}{2})\).

    \item \(\left(\sqrt{3}, -1\right)\)
    
    \textbf{Solution:}
    
    Take \(x = \sqrt{3}\), \(y = -1\). From \crefrange{eqn:polar-pythagorean}{eqn:polar-theta},
    \[
    r = \sqrt{(\sqrt{3})^2 + (-1)^2} = 2
    \]
    \(\theta\) lies in the \(4^{th}\) quadrant. So,
    \[
    \theta = \tan^{-1}\left(\frac{-1}{\sqrt{3}}\right) = -\frac{\pi}{6}
    \]
    
    \(\therefore\) The required point is \(_p(2, -\frac{\pi}{6})\).

    \item \(\left(-\frac{\sqrt{2}}{2}, \frac{\sqrt{2}}{2}\right)\)
    
    \textbf{Solution:}
    
    Take \(x = -\frac{\sqrt{2}}{2}\), \(y = \frac{\sqrt{2}}{2}\). From \crefrange{eqn:polar-pythagorean}{eqn:polar-theta},
    \[
    r = \sqrt{\left(-\frac{\sqrt{2}}{2}\right)^2 + \left(\frac{\sqrt{2}}{2}\right)^2} = 1
    \]
    \(\theta\) lies in the \(2^{nd}\) quadrant. So,
    \[
    \theta = \tan^{-1}(-1) + \pi = \frac{3\pi}{4}
    \]
    
    \(\therefore\) The required point is \(_p(1, \frac{3\pi}{4})\).

    \item \(\left(-2, 2\sqrt{3}\right)\)
    
    \textbf{Solution:}
    
    Take \(x = -2\), \(y = 2\sqrt{3}\). From \crefrange{eqn:polar-pythagorean}{eqn:polar-theta},
    \[
    r = \sqrt{(-2)^2 + (2\sqrt{3})^2} = 4
    \]
    \(\theta\) lies in the \(2^{nd}\) quadrant. So,
    \[
    \theta = \tan^{-1}\left(-\sqrt{3}\right) + \pi = \frac{2\pi}{3}
    \]
    
    \(\therefore\) The required point is \(_p(4, \frac{2\pi}{3})\).

    \item \(\left(\frac{5\sqrt{3}}{2}, -\frac{5}{2}\right)\)
    
    \textbf{Solution:}
    
    Take \(x = \frac{5\sqrt{3}}{2}\), \(y = -\frac{5}{2}\). From \crefrange{eqn:polar-pythagorean}{eqn:polar-theta},
    \[
    r = \sqrt{\left(\frac{5\sqrt{3}}{2}\right)^2 + \left(-\frac{5}{2}\right)^2} = 5
    \]
    \(\theta\) lies in the \(4^{th}\) quadrant. So,
    \[
    \theta = \tan^{-1}\left(-\frac{1}{\sqrt{3}}\right) = -\frac{\pi}{6}
    \]
    
    \(\therefore\) The required point is \(_p(5, -\frac{\pi}{6})\).

    \item \(\left(3\sqrt{2}, -3\sqrt{2}\right)\)
    
    \textbf{Solution:}
    
    Take \(x = 3\sqrt{2}\), \(y = -3\sqrt{2}\). From \crefrange{eqn:polar-pythagorean}{eqn:polar-theta},
    \[
    r = \sqrt{(3\sqrt{2})^2 + (-3\sqrt{2})^2} = 6
    \]
    \(\theta\) lies in the \(4^{th}\) quadrant. So,
    \[
    \theta = \tan^{-1}(-1) = -\frac{\pi}{4}
    \]
    
    \(\therefore\) The required point is \(_p(6, -\frac{\pi}{4})\).
\end{enumerate}

\subsection*{Convert Polar Coordinates to Rectangular Form}

\begin{enumerate}
    \item \(\left( 3, \frac{\pi}{6} \right)\)
    
    \textbf{Solution:}
    
    Using \cref{eqn:cart-x} and \cref{eqn:cart-y}:
    \[
    x = 3\cos\left(\frac{\pi}{6}\right) = 3 \cdot \frac{\sqrt{3}}{2} = \frac{3\sqrt{3}}{2}
    \]
    \[
    y = 3\sin\left(\frac{\pi}{6}\right) = 3 \cdot \frac{1}{2} = \frac{3}{2}
    \]
    
    \(\therefore\) The rectangular coordinates are \(\left( \frac{3\sqrt{3}}{2}, \frac{3}{2} \right)\).
    
    \item \(\left( 5\sqrt{2}, -135^\circ \right)\)
    
    \textbf{Solution:}
    
    Convert \(-135^\circ\) to radians: \(-\frac{3\pi}{4}\).
    
    Using \cref{eqn:cart-x} and \cref{eqn:cart-y}:
    \[
    x = 5\sqrt{2}\cos\left(-\frac{3\pi}{4}\right) = 5\sqrt{2} \cdot \left(-\frac{\sqrt{2}}{2}\right) = -5
    \]
    \[
    y = 5\sqrt{2}\sin\left(-\frac{3\pi}{4}\right) = 5\sqrt{2} \cdot \left(-\frac{\sqrt{2}}{2}\right) = -5
    \]
    
    \(\therefore\) The rectangular coordinates are \((-5, -5)\).
    
    \item \((-4, 450^\circ)\)
    
    \textbf{Solution:}
    
    First find equivalent angle between \(0^\circ\) and \(360^\circ\):
    \(450^\circ - 360^\circ = 90^\circ\).
    
    Using \cref{eqn:cart-x} and \cref{eqn:cart-y}:
    \[
    x = -4\cos(90^\circ) = -4 \cdot 0 = 0
    \]
    \[
    y = -4\sin(90^\circ) = -4 \cdot 1 = -4
    \]
    
    \(\therefore\) The rectangular coordinates are \((0, -4)\).
    
    \item \(\left( -2, \frac{11\pi}{6} \right)\)
    
    \textbf{Solution:}
    
    Using \cref{eqn:cart-x} and \cref{eqn:cart-y}:
    \[
    x = -2\cos\left(\frac{11\pi}{6}\right) = -2 \cdot \frac{\sqrt{3}}{2} = -\sqrt{3}
    \]
    \[
    y = -2\sin\left(\frac{11\pi}{6}\right) = -2 \cdot \left(-\frac{1}{2}\right) = 1
    \]
    
    \(\therefore\) The rectangular coordinates are \((-\sqrt{3}, 1)\).
    
    \item \(\left( -5, -\frac{5\pi}{6} \right)\)
    
    \textbf{Solution:}
    
    Using \cref{eqn:cart-x} and \cref{eqn:cart-y}:
    \[
    x = -5\cos\left(-\frac{5\pi}{6}\right) = -5 \cdot \left(-\frac{\sqrt{3}}{2}\right) = \frac{5\sqrt{3}}{2}
    \]
    \[
    y = -5\sin\left(-\frac{5\pi}{6}\right) = -5 \cdot \frac{1}{2} = -\frac{5}{2}
    \]
    
    \(\therefore\) The rectangular coordinates are \(\left( \frac{5\sqrt{3}}{2}, -\frac{5}{2} \right)\).
    
    \item \((3, -870^\circ)\)
    
    \textbf{Solution:}
    
    First find equivalent positive angle:
    \(-870^\circ + 3 \times 360^\circ = 210^\circ\).
    
    Using \cref{eqn:cart-x} and \cref{eqn:cart-y}:
    \[
    x = 3\cos(210^\circ) = 3 \cdot \left(-\frac{\sqrt{3}}{2}\right) = -\frac{3\sqrt{3}}{2}
    \]
    \[
    y = 3\sin(210^\circ) = 3 \cdot \left(-\frac{1}{2}\right) = -\frac{3}{2}
    \]
    
    \(\therefore\) The rectangular coordinates are \(\left( -\frac{3\sqrt{3}}{2}, -\frac{3}{2} \right)\).
    
    \item \((-12, 0^\circ)\)
    
    \textbf{Solution:}
    
    Using \cref{eqn:cart-x} and \cref{eqn:cart-y}:
    \[
    x = -12\cos(0^\circ) = -12 \cdot 1 = -12
    \]
    \[
    y = -12\sin(0^\circ) = -12 \cdot 0 = 0
    \]
    
    \(\therefore\) The rectangular coordinates are \((-12, 0)\).
    
    \item \(\left( \sqrt{6}, \frac{\pi}{3} \right)\)
    
    \textbf{Solution:}
    
    Using \cref{eqn:cart-x} and \cref{eqn:cart-y}:
    \[
    x = \sqrt{6}\cos\left(\frac{\pi}{3}\right) = \sqrt{6} \cdot \frac{1}{2} = \frac{\sqrt{6}}{2}
    \]
    \[
    y = \sqrt{6}\sin\left(\frac{\pi}{3}\right) = \sqrt{6} \cdot \frac{\sqrt{3}}{2} = \frac{3\sqrt{2}}{2}
    \]
    
    \(\therefore\) The rectangular coordinates are \(\left( \frac{\sqrt{6}}{2}, \frac{3\sqrt{2}}{2} \right)\).
\end{enumerate}


  
\end{document}

